% Options for packages loaded elsewhere
% Options for packages loaded elsewhere
\PassOptionsToPackage{unicode}{hyperref}
\PassOptionsToPackage{hyphens}{url}
\PassOptionsToPackage{dvipsnames,svgnames,x11names}{xcolor}
%
\documentclass[
  11pt,
  letterpaper,
]{article}
\usepackage{xcolor}
\usepackage[top=0.5in,left=1in,right=1in,bottom=0.75in]{geometry}
\usepackage{amsmath,amssymb}
\setcounter{secnumdepth}{-\maxdimen} % remove section numbering
\usepackage{iftex}
\ifPDFTeX
  \usepackage[T1]{fontenc}
  \usepackage[utf8]{inputenc}
  \usepackage{textcomp} % provide euro and other symbols
\else % if luatex or xetex
  \usepackage{unicode-math} % this also loads fontspec
  \defaultfontfeatures{Scale=MatchLowercase}
  \defaultfontfeatures[\rmfamily]{Ligatures=TeX,Scale=1}
\fi
\usepackage{lmodern}
\ifPDFTeX\else
  % xetex/luatex font selection
  \setmainfont[]{Arial}
\fi
% Use upquote if available, for straight quotes in verbatim environments
\IfFileExists{upquote.sty}{\usepackage{upquote}}{}
\IfFileExists{microtype.sty}{% use microtype if available
  \usepackage[]{microtype}
  \UseMicrotypeSet[protrusion]{basicmath} % disable protrusion for tt fonts
}{}
\makeatletter
\@ifundefined{KOMAClassName}{% if non-KOMA class
  \IfFileExists{parskip.sty}{%
    \usepackage{parskip}
  }{% else
    \setlength{\parindent}{0pt}
    \setlength{\parskip}{6pt plus 2pt minus 1pt}}
}{% if KOMA class
  \KOMAoptions{parskip=half}}
\makeatother
% Make \paragraph and \subparagraph free-standing
\makeatletter
\ifx\paragraph\undefined\else
  \let\oldparagraph\paragraph
  \renewcommand{\paragraph}{
    \@ifstar
      \xxxParagraphStar
      \xxxParagraphNoStar
  }
  \newcommand{\xxxParagraphStar}[1]{\oldparagraph*{#1}\mbox{}}
  \newcommand{\xxxParagraphNoStar}[1]{\oldparagraph{#1}\mbox{}}
\fi
\ifx\subparagraph\undefined\else
  \let\oldsubparagraph\subparagraph
  \renewcommand{\subparagraph}{
    \@ifstar
      \xxxSubParagraphStar
      \xxxSubParagraphNoStar
  }
  \newcommand{\xxxSubParagraphStar}[1]{\oldsubparagraph*{#1}\mbox{}}
  \newcommand{\xxxSubParagraphNoStar}[1]{\oldsubparagraph{#1}\mbox{}}
\fi
\makeatother

\usepackage{color}
\usepackage{fancyvrb}
\newcommand{\VerbBar}{|}
\newcommand{\VERB}{\Verb[commandchars=\\\{\}]}
\DefineVerbatimEnvironment{Highlighting}{Verbatim}{commandchars=\\\{\}}
% Add ',fontsize=\small' for more characters per line
\usepackage{framed}
\definecolor{shadecolor}{RGB}{241,243,245}
\newenvironment{Shaded}{\begin{snugshade}}{\end{snugshade}}
\newcommand{\AlertTok}[1]{\textcolor[rgb]{0.68,0.00,0.00}{#1}}
\newcommand{\AnnotationTok}[1]{\textcolor[rgb]{0.37,0.37,0.37}{#1}}
\newcommand{\AttributeTok}[1]{\textcolor[rgb]{0.40,0.45,0.13}{#1}}
\newcommand{\BaseNTok}[1]{\textcolor[rgb]{0.68,0.00,0.00}{#1}}
\newcommand{\BuiltInTok}[1]{\textcolor[rgb]{0.00,0.23,0.31}{#1}}
\newcommand{\CharTok}[1]{\textcolor[rgb]{0.13,0.47,0.30}{#1}}
\newcommand{\CommentTok}[1]{\textcolor[rgb]{0.37,0.37,0.37}{#1}}
\newcommand{\CommentVarTok}[1]{\textcolor[rgb]{0.37,0.37,0.37}{\textit{#1}}}
\newcommand{\ConstantTok}[1]{\textcolor[rgb]{0.56,0.35,0.01}{#1}}
\newcommand{\ControlFlowTok}[1]{\textcolor[rgb]{0.00,0.23,0.31}{\textbf{#1}}}
\newcommand{\DataTypeTok}[1]{\textcolor[rgb]{0.68,0.00,0.00}{#1}}
\newcommand{\DecValTok}[1]{\textcolor[rgb]{0.68,0.00,0.00}{#1}}
\newcommand{\DocumentationTok}[1]{\textcolor[rgb]{0.37,0.37,0.37}{\textit{#1}}}
\newcommand{\ErrorTok}[1]{\textcolor[rgb]{0.68,0.00,0.00}{#1}}
\newcommand{\ExtensionTok}[1]{\textcolor[rgb]{0.00,0.23,0.31}{#1}}
\newcommand{\FloatTok}[1]{\textcolor[rgb]{0.68,0.00,0.00}{#1}}
\newcommand{\FunctionTok}[1]{\textcolor[rgb]{0.28,0.35,0.67}{#1}}
\newcommand{\ImportTok}[1]{\textcolor[rgb]{0.00,0.46,0.62}{#1}}
\newcommand{\InformationTok}[1]{\textcolor[rgb]{0.37,0.37,0.37}{#1}}
\newcommand{\KeywordTok}[1]{\textcolor[rgb]{0.00,0.23,0.31}{\textbf{#1}}}
\newcommand{\NormalTok}[1]{\textcolor[rgb]{0.00,0.23,0.31}{#1}}
\newcommand{\OperatorTok}[1]{\textcolor[rgb]{0.37,0.37,0.37}{#1}}
\newcommand{\OtherTok}[1]{\textcolor[rgb]{0.00,0.23,0.31}{#1}}
\newcommand{\PreprocessorTok}[1]{\textcolor[rgb]{0.68,0.00,0.00}{#1}}
\newcommand{\RegionMarkerTok}[1]{\textcolor[rgb]{0.00,0.23,0.31}{#1}}
\newcommand{\SpecialCharTok}[1]{\textcolor[rgb]{0.37,0.37,0.37}{#1}}
\newcommand{\SpecialStringTok}[1]{\textcolor[rgb]{0.13,0.47,0.30}{#1}}
\newcommand{\StringTok}[1]{\textcolor[rgb]{0.13,0.47,0.30}{#1}}
\newcommand{\VariableTok}[1]{\textcolor[rgb]{0.07,0.07,0.07}{#1}}
\newcommand{\VerbatimStringTok}[1]{\textcolor[rgb]{0.13,0.47,0.30}{#1}}
\newcommand{\WarningTok}[1]{\textcolor[rgb]{0.37,0.37,0.37}{\textit{#1}}}

\usepackage{longtable,booktabs,array}
\newcounter{none} % for unnumbered tables
\usepackage{calc} % for calculating minipage widths
% Correct order of tables after \paragraph or \subparagraph
\usepackage{etoolbox}
\makeatletter
\patchcmd\longtable{\par}{\if@noskipsec\mbox{}\fi\par}{}{}
\makeatother
% Allow footnotes in longtable head/foot
\IfFileExists{footnotehyper.sty}{\usepackage{footnotehyper}}{\usepackage{footnote}}
\makesavenoteenv{longtable}
\usepackage{graphicx}
\makeatletter
\newsavebox\pandoc@box
\newcommand*\pandocbounded[1]{% scales image to fit in text height/width
  \sbox\pandoc@box{#1}%
  \Gscale@div\@tempa{\textheight}{\dimexpr\ht\pandoc@box+\dp\pandoc@box\relax}%
  \Gscale@div\@tempb{\linewidth}{\wd\pandoc@box}%
  \ifdim\@tempb\p@<\@tempa\p@\let\@tempa\@tempb\fi% select the smaller of both
  \ifdim\@tempa\p@<\p@\scalebox{\@tempa}{\usebox\pandoc@box}%
  \else\usebox{\pandoc@box}%
  \fi%
}
% Set default figure placement to htbp
\def\fps@figure{htbp}
\makeatother





\setlength{\emergencystretch}{3em} % prevent overfull lines

\providecommand{\tightlist}{%
  \setlength{\itemsep}{0pt}\setlength{\parskip}{0pt}}



 


\makeatletter
\@ifpackageloaded{caption}{}{\usepackage{caption}}
\AtBeginDocument{%
\ifdefined\contentsname
  \renewcommand*\contentsname{Table of contents}
\else
  \newcommand\contentsname{Table of contents}
\fi
\ifdefined\listfigurename
  \renewcommand*\listfigurename{List of Figures}
\else
  \newcommand\listfigurename{List of Figures}
\fi
\ifdefined\listtablename
  \renewcommand*\listtablename{List of Tables}
\else
  \newcommand\listtablename{List of Tables}
\fi
\ifdefined\figurename
  \renewcommand*\figurename{Figure}
\else
  \newcommand\figurename{Figure}
\fi
\ifdefined\tablename
  \renewcommand*\tablename{Table}
\else
  \newcommand\tablename{Table}
\fi
}
\@ifpackageloaded{float}{}{\usepackage{float}}
\floatstyle{ruled}
\@ifundefined{c@chapter}{\newfloat{codelisting}{h}{lop}}{\newfloat{codelisting}{h}{lop}[chapter]}
\floatname{codelisting}{Listing}
\newcommand*\listoflistings{\listof{codelisting}{List of Listings}}
\makeatother
\makeatletter
\makeatother
\makeatletter
\@ifpackageloaded{caption}{}{\usepackage{caption}}
\@ifpackageloaded{subcaption}{}{\usepackage{subcaption}}
\makeatother
\usepackage{bookmark}
\IfFileExists{xurl.sty}{\usepackage{xurl}}{} % add URL line breaks if available
\urlstyle{same}
\hypersetup{
  pdftitle={Research Memo - Apprehension Type Engineering},
  pdfauthor={Andrew K. Sanford},
  colorlinks=true,
  linkcolor={blue},
  filecolor={Maroon},
  citecolor={Blue},
  urlcolor={Blue},
  pdfcreator={LaTeX via pandoc}}


\title{Research Memo - Apprehension Type Engineering}
\author{Andrew K. Sanford}
\date{2026-07-30}
\begin{document}
\maketitle


\textbf{Re:} {[}How is Dominant Apprehension Type determined?{]}

Dominant Apprehension Type is an engineered outcome variable that
identifies the greatest of two composite scores: Group Discussion
Apprehension \& Solo Conversation Apprehension.

\begin{Shaded}
\begin{Highlighting}[]
\CommentTok{\# Identify forward and reverse{-}coded Q\textquotesingle{}s (Higher scores = more apprehensive)}
\NormalTok{forward\_code }\OtherTok{\textless{}{-}} \FunctionTok{c}\NormalTok{(}
  \StringTok{"q1\_dislike\_groupdiscussion"}\NormalTok{, }\StringTok{"q3\_nervous\_groupdiscussion"}\NormalTok{,}
  \StringTok{"q5\_tense\_newppl\_groupdiscussion"}\NormalTok{, }\StringTok{"q13\_nervous\_newppl\_convo"}\NormalTok{,}
  \StringTok{"q15\_nervous\_convo"}\NormalTok{, }\StringTok{"q18\_afraid\_convo"}
\NormalTok{  )}
\NormalTok{reverse\_code }\OtherTok{\textless{}{-}} \FunctionTok{c}\NormalTok{(}
  \StringTok{"q2\_comfort\_groupdiscussion"}\NormalTok{, }\StringTok{"q4\_like\_groupdiscussion"}\NormalTok{,}
  \StringTok{"q6\_relaxed\_groupdiscussion"}\NormalTok{, }\StringTok{"q14\_nofear\_speaking\_convo"}\NormalTok{,}
  \StringTok{"q16\_relaxed\_convo"}\NormalTok{, }\StringTok{"q17\_relaxed\_newppl\_convo"}
\NormalTok{  )}
\CommentTok{\# Set numeric values for Likert scales}
\NormalTok{likert\_forward }\OtherTok{\textless{}{-}} \FunctionTok{c}\NormalTok{(}
  \StringTok{"Strongly disagree"}          \OtherTok{=} \DecValTok{1}\NormalTok{,}
  \StringTok{"Somewhat disagree"}          \OtherTok{=} \DecValTok{2}\NormalTok{,}
  \StringTok{"Neither agree nor disagree"} \OtherTok{=} \DecValTok{3}\NormalTok{,}
  \StringTok{"Somewhat agree"}             \OtherTok{=} \DecValTok{4}\NormalTok{,}
  \StringTok{"Strongly agree"}             \OtherTok{=} \DecValTok{5}
\NormalTok{)}
\NormalTok{likert\_reverse }\OtherTok{\textless{}{-}} \FunctionTok{c}\NormalTok{(}
  \StringTok{"Strongly disagree"}          \OtherTok{=} \DecValTok{5}\NormalTok{,}
  \StringTok{"Somewhat disagree"}          \OtherTok{=} \DecValTok{4}\NormalTok{,}
  \StringTok{"Neither agree nor disagree"} \OtherTok{=} \DecValTok{3}\NormalTok{,}
  \StringTok{"Somewhat agree"}             \OtherTok{=} \DecValTok{2}\NormalTok{,}
  \StringTok{"Strongly agree"}             \OtherTok{=} \DecValTok{1}
\NormalTok{)}
\CommentTok{\# Convert to numeric}
\NormalTok{survey }\OtherTok{\textless{}{-}}\NormalTok{ survey }\SpecialCharTok{|\textgreater{}}
  \FunctionTok{mutate}\NormalTok{(}
    \FunctionTok{across}\NormalTok{(}\FunctionTok{all\_of}\NormalTok{(forward\_code), }\SpecialCharTok{\textasciitilde{}}\NormalTok{ likert\_forward[.x]),}
    \FunctionTok{across}\NormalTok{(}\FunctionTok{all\_of}\NormalTok{(reverse\_code), }\SpecialCharTok{\textasciitilde{}}\NormalTok{ likert\_reverse[.x])}
\NormalTok{  )}
\CommentTok{\# Identify group discussion vs. solo conversation Q\textquotesingle{}s}
\NormalTok{group\_discussion }\OtherTok{\textless{}{-}} \FunctionTok{c}\NormalTok{(}
  \StringTok{"q1\_dislike\_groupdiscussion"}\NormalTok{, }\StringTok{"q2\_comfort\_groupdiscussion"}\NormalTok{,}
  \StringTok{"q3\_nervous\_groupdiscussion"}\NormalTok{, }\StringTok{"q4\_like\_groupdiscussion"}\NormalTok{,}
  \StringTok{"q6\_relaxed\_groupdiscussion"}
\NormalTok{  )}
\NormalTok{solo\_conversation }\OtherTok{\textless{}{-}} \FunctionTok{c}\NormalTok{(}
  \StringTok{"q13\_nervous\_newppl\_convo"}\NormalTok{, }\StringTok{"q14\_nofear\_speaking\_convo"}\NormalTok{,}
  \StringTok{"q15\_nervous\_convo"}\NormalTok{,        }\StringTok{"q16\_relaxed\_convo"}\NormalTok{,}
  \StringTok{"q17\_relaxed\_newppl\_convo"}\NormalTok{, }\StringTok{"q18\_afraid\_convo"}
\NormalTok{  )}
\CommentTok{\# Calculate biggest fear}
\NormalTok{survey }\OtherTok{\textless{}{-}}\NormalTok{ survey }\SpecialCharTok{|\textgreater{}} 
  \FunctionTok{mutate}\NormalTok{(}
    
\CommentTok{\#   Average of group apprehension Q\textquotesingle{}s}
    \AttributeTok{group\_apprehension =} \FunctionTok{rowMeans}\NormalTok{(}
      \FunctionTok{pick}\NormalTok{(}\FunctionTok{all\_of}\NormalTok{(group\_discussion)), }\AttributeTok{na.rm =} \ConstantTok{TRUE}\NormalTok{),}
    
\CommentTok{\#   Average of solo apprehension Q\textquotesingle{}s}
    \AttributeTok{solo\_apprehension =} \FunctionTok{rowMeans}\NormalTok{(}
      \FunctionTok{pick}\NormalTok{(}\FunctionTok{all\_of}\NormalTok{(solo\_conversation)), }\AttributeTok{na.rm =} \ConstantTok{TRUE}\NormalTok{),}

\CommentTok{\#   Assign greatest of the two values (equal becomes NA)}
    \AttributeTok{biggest\_fear =} \FunctionTok{case\_when}\NormalTok{(}
\NormalTok{      group\_apprehension }\SpecialCharTok{\textgreater{}}\NormalTok{ solo\_apprehension }\SpecialCharTok{\textasciitilde{}} \StringTok{"group\_apprehension"}\NormalTok{,}
\NormalTok{      solo\_apprehension }\SpecialCharTok{\textgreater{}}\NormalTok{ group\_apprehension }\SpecialCharTok{\textasciitilde{}} \StringTok{"solo\_apprehension"}\NormalTok{,}
      \AttributeTok{.default =} \StringTok{"equal"}\NormalTok{)}
\NormalTok{  )}
\end{Highlighting}
\end{Shaded}

After numericizing Likert scale values from a participants PRCA
responses, relevant questions were sorted into these two groups,
depending on which type of apprehension each question targeted. Each
group then had the corresponding numeric Likert scale values averaged to
calculate the composite scores for `average group discussion
apprehension' and `average solo conversation apprehension'. Finally, the
Dominant Apprehension Type takes the largest of these two values and
codes for the corresponding group, or indicates if both scores are
equal.

\emph{NOTE: This code for engineering `Dominant Transportation Method'
occurs in the clean\_data.qmd script.}




\end{document}
